

  \ch{Appendix A: Survey Instrument and Research Question Mapping}\label{appendix:survey-mapping}

  This appendix reproduces the survey instrument as administered and provides traceability between survey questions and research questions, following the systems engineering lifecycle approach outlined in the dissertation methodology. The survey structure addresses three core dimensions (\textbf{Awareness}, \textbf{Applicability}, and \textbf{Perceived Value}) as they pertain to Digital Engineering and its four pillars: Model-Based Systems Engineering, the Digital Thread, Digital Twin, and Product Lifecycle Management.

  \section{Survey Instrument as Administered}

  Participants first viewed an informed-consent notice. It stated that the survey was part of a doctoral research project at Dakota State University; that participation was anonymous, voluntary, and limited to respondents aged eighteen or older; that no personally identifiable information was collected; that risks were minimal and no compensation was offered; and that completion and submission constituted informed consent. The instrument used two five-point response scales, presented in Table~\ref{tab:scales}: a familiarity scale for the Section~1 awareness item and an agreement scale for the Likert items in Sections~2 through~5.

  \begin{table}[H]
    \centering
    \caption{The two five-point response scales used in the survey instrument.}\label{tab:scales}
    % ALT: Maps each of the five scale points to its label on the Section 1 familiarity scale and the Sections 2 to 5 agreement scale.
    \begin{tabular}{@{}cll@{}}
      \toprule
      \textbf{Point} & \textbf{Familiarity scale (Section 1)} & \textbf{Agreement scale (Sections 2--5)} \\
      \midrule
      1 & Not at all familiar & Strongly disagree \\
      2 & Slightly familiar   & Disagree \\
      3 & Moderately familiar & Neither agree nor disagree \\
      4 & Very familiar       & Agree \\
      5 & Extremely familiar  & Strongly agree \\
      \bottomrule
    \end{tabular}
  \end{table}

  The thirty items, by section, were as follows.

  \paragraph{Section 1: Awareness and Familiarity.}
  \begin{enumerate}[leftmargin=2.4em,itemsep=1pt]
    \item[Q1.1] Please rate your level of familiarity with Digital Engineering concepts and practices. \emph{(Familiarity scale)}
    \item[Q1.2] Have you encountered Digital Engineering methodologies, frameworks, or tools in your professional work within the past two years? \emph{(Yes / No)}
  \end{enumerate}

  \paragraph{Section 2: Understanding of Capabilities.} \emph{(Agreement scale)}
  \begin{enumerate}[leftmargin=2.4em,itemsep=1pt]
    \item[Q2.1] Digital Engineering includes model-based systems engineering approaches that can improve development processes.
    \item[Q2.2] Digital Engineering can enable digital twin development and virtual prototyping for Information Technology systems.
    \item[Q2.3] Digital Engineering supports continuous integration and data-driven decision-making in technology development.
    \item[Q2.4] Digital Engineering enables digital twin technology that can simulate security scenarios and test defensive measures without impacting production systems.
    \item[Q2.5] Digital Engineering supports continuous security validation and data-driven threat analysis throughout the development lifecycle.
    \item[Q2.6] Digital Engineering can improve security control implementation through automated compliance checking and verification.
  \end{enumerate}

  \paragraph{Section 3: Applicability.} \emph{(Agreement scale)}
  \begin{enumerate}[leftmargin=2.4em,itemsep=1pt]
    \item[Q3.1] Digital Engineering methodologies have relevant applications within the Information Technology sector.
    \item[Q3.2] Digital Engineering methodologies have relevant applications for addressing Information Assurance challenges.
    \item[Q3.3] The ability to utilize digital twins to test changes against accurate replicas of production environments would provide value to my organization.
    \item[Q3.4] The use of digital models to map and document an organization's environment and configurations would provide value to my organization.
    \item[Q3.5] The use of digital lifecycle management to meet compliance and service delivery requirements would provide value to my organization.
    \item[Q3.6] My organization faces regulatory or compliance requirements that could benefit from Digital Engineering approaches.
  \end{enumerate}

  \paragraph{Section 4: Value for Information Technology.} \emph{(Agreement scale; Q4.5 Yes / No / Unsure)}
  \begin{enumerate}[leftmargin=2.4em,itemsep=1pt]
    \item[Q4.1] Digital Engineering could deliver meaningful value to my organization's Information Technology processes.
    \item[Q4.2] Digital Engineering could reduce development cycle time in my organization.
    \item[Q4.3] Digital Engineering could improve product quality and reduce defects in my organization.
    \item[Q4.4] Digital Engineering could improve collaboration effectiveness across development teams in my organization.
    \item[Q4.5] My organization would be willing to invest in Digital Engineering capabilities if clear return on investment could be demonstrated.
  \end{enumerate}

  \paragraph{Section 5: Value for Information Assurance and Cybersecurity.} \emph{(Agreement scale; Q5.7 Yes / No / Unsure)}
  \begin{enumerate}[leftmargin=2.4em,itemsep=1pt]
    \item[Q5.1] Digital Engineering could deliver meaningful value to my organization's Information Assurance and cybersecurity operations.
    \item[Q5.2] Digital Engineering could reduce the time required to identify and remediate security vulnerabilities in my organization.
    \item[Q5.3] Digital Engineering could improve security posture and reduce successful cyber incidents in my organization.
    \item[Q5.4] Digital Engineering could enhance threat modeling and risk assessment capabilities in my organization.
    \item[Q5.5] Digital Engineering could improve collaboration between security teams, development teams, and operations teams in my organization.
    \item[Q5.6] Digital Engineering could help my organization achieve better compliance with security frameworks and regulatory requirements.
    \item[Q5.7] My organization would be willing to invest in Digital Engineering capabilities for cybersecurity purposes if clear return on investment could be demonstrated.
  \end{enumerate}

  \paragraph{Section 6: Interest and Demographics.}
  \begin{enumerate}[leftmargin=2.4em,itemsep=1pt]
    \item[Q6.1] Would you be interested in learning more about Digital Engineering applications for Information Assurance, Cybersecurity, or Information Technology in your industry? \emph{(Yes / No)}
    \item[Q6.2] Would you recommend that your organization explore Digital Engineering adoption for improving security operations? \emph{(Yes / No)}
    \item[Q6.3] Please indicate your field of practice. \emph{(Information Technology / Information Assurance or Cyber Security / Engineering / Other)}
    \item[Q6.4] Please indicate your level of experience in your field of practice. \emph{(1--5 / 6--10 / 11--15 / 16+ years)}
  \end{enumerate}

  Consistent with the design described in Chapter~3, the Section~2 through Section~5 stems carried only minimal functional context for each capability; respondents did not see a separate definition of Digital Engineering, which bears on the interpretation of the high endorsement reported in Chapter~4.

  \paragraph{Instrument Limitations.} Because the instrument is reproduced as administered, several design limitations bear on interpretation rather than on revision. Several stems bundle related capabilities in a single item (e.g., Q2.3, Q2.5, Q3.5, Q5.2, Q5.3, Q5.6), so item-level agreement cannot be attributed to any single clause. All Likert items carry positive keying, which leaves residual acquiescence risk; the straight-lining screen in Appendix~B diagnoses that risk. The phrasing of Q3.3 (``accurate replicas'') presupposes digital-twin fidelity rather than asking respondents to evaluate it. Q4.5 and Q5.7 pose hypothetical organizational investment questions that individual respondents answered on their organizations' behalf. Finally, the demographic items offer no ``prefer not to answer'' option, and the experience item (Q6.4) provides no category below one year.

  \section{Recruitment Announcement as Posted}\label{sec:recruitment-announcement}

  Recruitment proceeded primarily through a LinkedIn post, as described in Chapter~3. That announcement was deliberately descriptive rather than persuasive. It named the study topic and the eligible population without asserting that Digital Engineering delivers value, so that its wording would not prime respondents toward agreement. Below, the post appears verbatim as administered.\footnote{The announcement stated a ``31 question survey.'' One item was subsequently identified as a duplicate and removed before administration, so the instrument as administered contained thirty items, listed in the Survey Instrument as Administered section above. The recruitment text is reproduced here with its original count.}

  \begin{quote}
    \textbf{Participants Wanted for Research Study}\\
    We are conducting a research study about Digital Engineering, Information Assurance, and Information Technology Service Management, and are inviting individuals to participate.

    \textbf{Who can participate?}
    \begin{itemize}[leftmargin=1.6em,nosep]
      \item IT, Information Assurance, and Cybersecurity Professionals
    \end{itemize}

    \textbf{What will you do?}
    \begin{itemize}[leftmargin=1.6em,nosep]
      \item Complete a 31 question survey
      \item Estimated time: 10 minutes
    \end{itemize}

    \textbf{No Compensation}

    \textbf{Voluntary Participation}\\
    Your participation is completely voluntary.

    \textbf{Interested?}\\
    Click here to learn more or participate: \texttt{https://lnkd.in/guNhiyky}\\
    For questions, contact: John Bonar, \texttt{john.bonar@trojans.dsu.edu}
  \end{quote}

  The paid LinkedIn promotion described in Chapter~3 boosted the visibility of this same post. It altered reach only, and changed neither the wording nor the incentive structure. Because the survey was administered anonymously through Microsoft Forms, the platform did not attribute individual responses to specific referral channels, so channel-level response counts are not available. Chapter~3 characterizes the sector composition of the realized sample, and Chapter~5 addresses its implications for generalizability.

  \section{Research Questions}

  \begin{table}[H]
    \centering
    \caption{Research questions.}\label{tab:research-questions}
    % ALT: Lists the three research questions, pairing each RQ identifier (RQ1, RQ2, RQ3) with its full research-question text.
    \begin{tabular}{@{}cp{5.5in}@{}}
      \toprule
      \textbf{RQ} & \textbf{Research Question} \\
      \midrule
      RQ1 & To what extent are IT and Information Assurance professionals aware of Digital Engineering capabilities, including MBSE, the Digital Thread, digital twin technologies, and PLM principles? \\
      RQ2 & Do IT and Information Assurance professionals perceive Digital Engineering capabilities as potentially valuable or important for their work in IA, security compliance, and IT service delivery? \\
      RQ3 & Do IT and Information Assurance professionals believe DE practices would help them perform their jobs, meet compliance requirements, or enhance organizational capabilities? \\
      \bottomrule
    \end{tabular}
  \end{table}
  \clearpage
  \iffalse % Orphan figure (never referenced in text) commented out per author request
  \begin{landscape}
    \begin{figure}
      \centering
      \includegraphics[width=\linewidth, alt={SysML requirement diagram showing the three research questions as requirement elements, each with an identifier and text, linked by Satisfy dependencies to the survey-question Blocks that measure it; each Block carries a Likert, Binary, or Ternary response operation.}]{figures/ResearchQuestionTaxonomy.pdf}
      \caption{Research Question Traceability Taxonomy}
      \label{fig:SurveyTaxonomyAppendixA}
    \end{figure}
  \end{landscape}
  \fi
  \section{Survey Structure Aligned to Core Dimensions}

  \begin{table}[H]
    \centering
    \caption{Survey section structure.}
    \label{tab:survey-structure}
    % ALT: For each of the six survey sections, gives its core dimension, primary research question, and a short description.
    \begin{tabular}{@{}lllp{2.5in}@{}}
      \toprule
      \textbf{Section} & \textbf{Core Dimension} & \textbf{Primary RQ} & \textbf{Description} \\
      \midrule
      Section 1 & Awareness & RQ1 & Baseline familiarity and professional exposure \\
      Section 2 & Awareness & RQ1 & Understanding of specific DE capabilities/pillars \\
      Section 3 & Applicability & RQ2, RQ3 & Perceived relevance to IT and IA domains \\
      Section 4 & Perceived Value & RQ2, RQ3 & Value assessment for IT operations \\
      Section 5 & Perceived Value & RQ2, RQ3 & Value assessment for IA/Cybersecurity operations \\
      Section 6 & Demographics & Supporting & Professional field and experience level \\
      \bottomrule
    \end{tabular}
  \end{table}

  \iffalse % Orphan figure (never referenced in text) commented out per author request
  \begin{landscape}
    \begin{figure}
      \centering
      \includegraphics[width=.9\linewidth, alt={Taxonomy of the survey-question elements organized by the six survey sections.}]{figures/OverviewDiagramOfPackageSurveyQuestions.pdf}
      \caption{Taxonomy of Survey Questions}
      \label{fig:SurveyTaxonomy}
    \end{figure}
  \end{landscape}
  \fi
  \clearpage

  \section{Section 1: Awareness and Familiarity with Digital Engineering}

  \textbf{Core Dimension: Awareness}

  \textbf{Primary Research Question: RQ1}

  \begin{table}[H]
    \centering
    \caption{Section 1 question mapping.}
    \label{tab:section1-mapping}
    % ALT: Maps each Section 1 question (1.1, 1.2) to its question text, response format, DE pillar, and research question.
    \small
    \begin{tabular}{@{}cp{2.2in}p{1in}ll@{}}
      \toprule
      \textbf{Q\#} & \textbf{Question Text} & \textbf{Format} & \textbf{DE Pillar} & \textbf{RQ} \\
      \midrule
      1.1 & Please rate your level of familiarity with Digital Engineering concepts and practices. & 5-point Likert (Familiarity) & All/General & RQ1 \\
      1.2 & Have you encountered Digital Engineering methodologies, frameworks, or tools in your professional work within the past two years? & Binary (Yes/No) & All/General & RQ1 \\
      \bottomrule
    \end{tabular}
  \end{table}

  \subsection{Rationale}

  These questions establish baseline awareness metrics essential to all subsequent analysis. Question 1.1 measures self-assessed familiarity (theoretical awareness), while Question 1.2 measures practical professional exposure (applied awareness). Together they distinguish between those who have merely heard of Digital Engineering and those who have encountered it in the course of their professional duties.

  \section{Section 2: Understanding of Digital Engineering Capabilities}

  \textbf{Core Dimension: Awareness}

  \textbf{Primary Research Question: RQ1}

  \begin{table}[H]
    \centering
    \caption{Section 2 question mapping.}
    \label{tab:section2-mapping}
    % ALT: Maps each Section 2 question (2.1 to 2.6) to its question text, response format, DE pillar, and research question.
    \small
    \begin{tabular}{@{}cp{2.8in}p{0.9in}p{0.9in}l@{}}
      \toprule
      \textbf{Q\#} & \textbf{Question Text} & \textbf{Format} & \textbf{DE Pillar} & \textbf{RQ} \\
      \midrule
      2.1 & DE includes model-based systems engineering approaches that can improve development processes. & 5-point Likert (Agreement) & MBSE & RQ1 \\
      2.2 & DE can enable digital twin development and virtual prototyping for IT systems. & 5-point Likert (Agreement) & Digital Twin & RQ1 \\
      2.3 & DE supports continuous integration and data-driven decision-making in technology development. & 5-point Likert (Agreement) & Digital Thread & RQ1 \\
      2.4 & DE enables digital twin technology that can simulate security scenarios and test defensive measures without impacting production systems. & 5-point Likert (Agreement) & Digital Twin (Security) & RQ1 \\
      2.5 & DE supports continuous security validation and data-driven threat analysis throughout the development lifecycle. & 5-point Likert (Agreement) & Digital Thread (Security) & RQ1 \\
      2.6 & DE can improve security control implementation through automated compliance checking and verification. & 5-point Likert (Agreement) & PLM / Traceability & RQ1 \\
      \bottomrule
    \end{tabular}
  \end{table}

  \subsection{Rationale}

  This section probes understanding of specific Digital Engineering pillars applied to IT and security contexts. Questions address all four pillars: MBSE (Q2.1), Digital Twin (Q2.2, Q2.4), the Digital Thread (Q2.3, Q2.5), and PLM with its authoritative traceability capabilities (Q2.6). Including both general IT applications (Q2.1-Q2.3) and security-specific applications (Q2.4-Q2.6) enables meaningful comparison across professional domains.

  \section{Section 3: Applicability of Digital Engineering}

  \textbf{Core Dimension: Applicability}

  \textbf{Primary Research Questions: RQ2, RQ3}

  \begin{table}[H]
    \centering
    \caption{Section 3 question mapping.}
    \label{tab:section3-mapping}
    % ALT: Maps each Section 3 question (3.1 to 3.6) to its question text, response format, DE pillar, and research question.
    \small
    \begin{tabular}{@{}cp{2.2in}p{0.9in}p{0.8in}l@{}}
      \toprule
      \textbf{Q\#} & \textbf{Question Text} & \textbf{Format} & \textbf{DE Pillar} & \textbf{RQ} \\
      \midrule
      3.1 & DE methodologies have relevant applications within the Information Technology sector. & 5-point Likert (Agreement) & All/General & RQ2 \\
      3.2 & DE methodologies have relevant applications for addressing Information Assurance challenges. & 5-point Likert (Agreement) & All/General & RQ2 \\
      3.3 & The ability to utilize digital twins to test changes against accurate replicas of production environments would provide value to my organization. & 5-point Likert (Agreement) & Digital Twin & RQ2, RQ3 \\
      3.4 & The use of digital models to map and document an organization's environment and configurations would provide value to my organization. & 5-point Likert (Agreement) & MBSE & RQ2, RQ3 \\
      3.5 & The use of digital lifecycle management to meet compliance and service delivery requirements would provide value to my organization. & 5-point Likert (Agreement) & PLM & RQ2, RQ3 \\
      3.6 & My organization faces regulatory or compliance requirements that could benefit from Digital Engineering approaches. & 5-point Likert (Agreement) & All/General & RQ2, RQ3 \\
      \bottomrule
    \end{tabular}
  \end{table}

  \subsection{Rationale}

  This section bridges awareness and value by examining whether respondents perceive Digital Engineering as applicable to their professional contexts. Questions 3.1 and 3.2 assess domain-level applicability (IT versus IA), while Questions 3.3 through 3.5 assess pillar-specific organizational value (Digital Twin, MBSE, PLM). Question 3.6 identifies compliance-driven need, which is central to the research focus on Information Assurance and regulatory compliance.

  \section{Section 4: Value Assessment for Information Technology}

  \textbf{Core Dimension: Perceived Value (IT Domain)}

  \textbf{Primary Research Questions: RQ2, RQ3}

  \begin{table}[H]
    \centering
    \caption{Section 4 question mapping.}
    \label{tab:section4-mapping}
    % ALT: Maps each Section 4 question (4.1 to 4.5) to its question text, response format, value category, and research question.
    \small
    \begin{tabular}{@{}cp{2in}p{0.9in}p{1in}l@{}}
      \toprule
      \textbf{Q\#} & \textbf{Question Text} & \textbf{Format} & \textbf{Value Category} & \textbf{RQ} \\
      \midrule
      4.1 & DE could deliver meaningful value to my organization's Information Technology processes. & 5-point Likert (Agreement) & Overall IT Value & RQ2, RQ3 \\
      4.2 & DE could reduce development cycle time in my organization. & 5-point Likert (Agreement) & Efficiency Benefit & RQ3 \\
      4.3 & DE could improve product quality and reduce defects in my organization. & 5-point Likert (Agreement) & Quality Benefit & RQ3 \\
      4.4 & DE could improve collaboration effectiveness across development teams in my organization. & 5-point Likert (Agreement) & Collaboration Benefit & RQ3 \\
      4.5 & My organization would be willing to invest in DE capabilities if clear ROI could be demonstrated. & Ternary (Yes/No/Unsure) & Investment Willingness & RQ2 \\
      \bottomrule
    \end{tabular}
  \end{table}

  \subsection{Rationale}

  This section measures IT-specific value perceptions. Question 4.1 provides an overall IT value assessment. Questions 4.2 through 4.4 examine specific operational benefits (efficiency, quality, and collaboration) that directly relate to job performance as addressed by RQ3. Question 4.5 measures organizational adoption interest, indicating whether perceived value rises to a level sufficient to warrant investment.

  \section{Section 5: Value Assessment for Information Assurance}

  \textbf{Core Dimension: Perceived Value (Information Assurance Domain)}

  \textbf{Primary Research Questions: RQ2, RQ3}

  \begin{table}[H]
    \centering
    \caption{Section 5 question mapping.}
    \label{tab:section5-mapping}
    % ALT: Maps each Section 5 question (5.1 to 5.7) to its question text, response format, value category, and research question.
    \small
    \begin{tabular}{@{}cp{2.5in}p{1.0in}p{1.0in}l@{}}
      \toprule
      \textbf{Q\#} & \textbf{Question Text} & \textbf{Format} & \textbf{Value Category} & \textbf{RQ} \\
      \midrule
      5.1 & DE could deliver meaningful value to my organization's Information Assurance and cybersecurity operations. & 5-point Likert (Agreement) & Overall Security Value & RQ2, RQ3 \\
      5.2 & DE could reduce the time required to identify and remediate security vulnerabilities in my organization. & 5-point Likert (Agreement) & Vulnerability Mgmt Benefit & RQ3 \\
      5.3 & DE could improve security posture and reduce successful cyber incidents in my organization. & 5-point Likert (Agreement) & Security Posture Benefit & RQ3 \\
      5.4 & DE could enhance threat modeling and risk assessment capabilities in my organization. & 5-point Likert (Agreement) & Threat Modeling Benefit & RQ3 \\
      5.5 & DE could improve collaboration between security teams, development teams, and operations teams in my organization. & 5-point Likert (Agreement) & Cross-Team Collaboration & RQ3 \\
      5.6 & DE could help my organization achieve better compliance with security frameworks and regulatory requirements. & 5-point Likert (Agreement) & Compliance Benefit & RQ3 \\
      5.7 & My organization would be willing to invest in DE capabilities for cybersecurity purposes if clear ROI could be demonstrated. & Ternary (Yes/No/Unsure) & Investment Willingness & RQ2 \\
      \bottomrule
    \end{tabular}
  \end{table}

  \subsection{Rationale}

  This section measures Information Assurance-specific value perceptions. Question 5.1 provides an overall security value assessment. Questions 5.2 through 5.6 examine specific security operational benefits directly related to job performance and compliance requirements as addressed by RQ3. The emphasis on compliance (Q5.6) directly addresses the dissertation's focus on Information Assurance compliance frameworks. Question 5.7 measures security-specific investment willingness.

  \section{Section 6: Interest and Demographic Information}

  \textbf{Core Dimension: Supporting/Demographics}

  \textbf{Purpose}

  Enable subgroup analysis and assess future research/adoption interest

  \begin{table}[H]
    \centering
    \caption{Section 6 question mapping.}
    \label{tab:section6-mapping}
    % ALT: Maps each Section 6 question (6.1 to 6.4) to its question text, response format, category, and research question.
    \small
    \begin{tabular}{@{}cp{2in}p{1.1in}p{1in}l@{}}
      \toprule
      \textbf{Q\#} & \textbf{Question Text} & \textbf{Format} & \textbf{Category} & \textbf{RQ} \\
      \midrule
      6.1 & Would you be interested in learning more about DE applications for Information Assurance, cybersecurity, or Information Technology in your industry? & Binary (Yes/No) & Learning Interest & Supporting \\
      6.2 & Would you recommend that your organization explore DE adoption for improving security operations? & Binary (Yes/No) & Recommendation & Supporting \\
      6.3 & Please indicate your field of practice. & Categorical (IT / IA or Cyber Security / Engineering / Other) & Professional Field & Demographics \\
      6.4 & Please indicate your level of experience in your field of practice. & Categorical (Experience ranges) & Experience Level & Demographics \\
      \bottomrule
    \end{tabular}
  \end{table}

  \subsection{Rationale}

  Questions 6.1 and 6.2 measure forward-looking interest that complements value perception, indicating whether positive perceptions translate into desire for learning and willingness to recommend organizational exploration. Questions 6.3 and 6.4 enable comparative analysis between IT and security professionals and across experience levels, addressing whether awareness and perceptions vary systematically by professional background.

  \section{Summary: Question Distribution by Research Question}

  \begin{table}[H]
    \centering
    \caption{Question distribution by research question.}
    \label{tab:question-distribution-rq}
    % ALT: For each research question (RQ1 to RQ3) and the demographics group, lists its primary questions, supporting questions, and total item count.
    \begin{tabular}{@{}lp{2in}p{1.2in}c@{}}
      \toprule
      \textbf{Research Question} & \textbf{Primary Questions} & \textbf{Supporting Questions} & \textbf{Total} \\
      \midrule
      RQ1 (Awareness) & Q1.1, Q1.2, Q2.1-Q2.6 & --- & 8 \\
      RQ2 (Perceived Value) & Q3.1-Q3.6, Q4.1, Q4.5, Q5.1, Q5.7 & Q6.1, Q6.2 & 12 \\
      RQ3 (Job/Compliance Help) & Q3.3-Q3.6, Q4.1-Q4.4, Q5.1-Q5.6 & --- & 14 \\
      Demographics & Q6.3, Q6.4 & --- & 2 \\
      \bottomrule
    \end{tabular}
  \end{table}

  Note: Several questions map to multiple research questions as they address both perceived value (RQ2) and anticipated benefits for job performance and compliance (RQ3).

  \section{Summary: Question Distribution by DE Pillar}

  \begin{table}[H]
    \centering
    \caption{Question distribution by Digital Engineering pillar.}
    \label{tab:question-distribution-pillar}
    % ALT: Lists the survey questions grouped under each Digital Engineering pillar (MBSE, Digital Twin, Digital Thread, PLM, and general/all pillars).
    \begin{tabular}{@{}lp{4in}@{}}
      \toprule
      \textbf{DE Pillar} & \textbf{Questions} \\
      \midrule
      MBSE & Q2.1, Q3.4 \\
      Digital Twin & Q2.2, Q2.4, Q3.3 \\
      Digital Thread & Q2.3, Q2.5 \\
      PLM & Q2.6, Q3.5 \\
      General/All Pillars & Q1.1, Q1.2, Q3.1, Q3.2, Q3.6, Q4.1-Q4.5, Q5.1-Q5.7, Q6.1-Q6.4 \\
      \bottomrule
    \end{tabular}
  \end{table}

  \section{Analysis Framework}

  \subsection{Primary Analysis for Each Research Question}

  \subsubsection{RQ1 Analysis (Awareness):}
  \begin{itemize}
    \item Mean familiarity score (Q1.1) with 95\% confidence interval
    \item Percentage with professional exposure (Q1.2) with 95\% confidence interval
    \item Mean agreement scores for capability understanding (Q2.1-Q2.6)
    \item Percentage indicating agreement (\(\geq\)4) with each capability statement
    \item Comparison of awareness levels across professional fields (Q6.3) and experience levels (Q6.4)
  \end{itemize}

  \subsubsection{RQ2 Analysis (Perceived Value):}
  \begin{itemize}
    \item Mean agreement scores for applicability and pillar-specific organizational value (Q3.1-Q3.6)
    \item Mean agreement scores for overall value (Q4.1, Q5.1)
    \item Percentage indicating investment willingness (Q4.5, Q5.7)
    \item Percentage indicating learning interest (Q6.1) and recommendation (Q6.2)
    \item Comparison of value perceptions across professional fields and experience levels
  \end{itemize}

  \subsubsection{RQ3 Analysis (Job/Compliance Help):}
  \begin{itemize}
    \item Mean agreement scores for specific benefits (Q4.2-Q4.4, Q5.2-Q5.6)
    \item Percentage agreeing with compliance benefits (Q3.5, Q3.6, Q5.6)
    \item Percentage agreeing with organizational capability benefits (Q3.3-Q3.5)
    \item Comparison of benefit perceptions across professional fields and experience levels
  \end{itemize}

  \subsection{Composite Scores}

  \begin{table}[H]
    \centering
    \caption{Composite score definitions.}
    \label{tab:composite-scores}
    % ALT: Defines each composite score by its constituent questions, item count, and interpretation.
    \setlength{\tabcolsep}{5pt}
    \begin{tabular}{@{}llcp{1.6in}@{}}
      \toprule
      \textbf{Composite} & \textbf{Questions} & \textbf{Items} & \textbf{Interpretation} \\
      \midrule
      Awareness Composite & Q1.1, Q2.1-Q2.6 & 7 & Higher = greater awareness/understanding \\
      IT Value Composite & Q3.1, Q3.3-Q3.5, Q4.1-Q4.4 & 8 & Higher = greater perceived IT value \\
      Security Value Composite & Q3.2, Q3.6, Q5.1-Q5.6 & 8 & Higher = greater perceived security value \\
      \bottomrule
    \end{tabular}
  \end{table}

  The Awareness composite averages the familiarity-anchored item Q1.1 with the agreement-anchored items Q2.1--Q2.6; the composite rests on the assumption that both item types, as five-point ordinal scales, are comparable, and Appendix~B notes the larger variance of Q1.1 relative to the Section~2 items. Cronbach's alpha assessed internal consistency (threshold $\alpha \geq 0.70$); Appendix~B reports the values.

  \section{Scale Reference}

  \subsection{Familiarity Scale (Q1.1)}

  \begin{table}[H]
    \centering
    \caption{Familiarity scale.}
    \label{tab:familiarity-scale}
    % ALT: Maps each familiarity-scale score (1 to 5) to its label and interpretation.
    \begin{tabular}{@{}cll@{}}
      \toprule
      \textbf{Score} & \textbf{Label} & \textbf{Interpretation} \\
      \midrule
      1 & Not at all familiar & No awareness \\
      2 & Slightly familiar & Minimal awareness \\
      3 & Moderately familiar & Basic awareness \\
      4 & Very familiar & Good awareness \\
      5 & Extremely familiar & Expert awareness \\
      \bottomrule
    \end{tabular}
  \end{table}

  \subsection{Agreement Scale (Q2.1--Q4.4 and Q5.1--Q5.6)}

  \begin{table}[H]
    \centering
    \caption{Agreement scale.}
    \label{tab:agreement-scale}
    % ALT: Maps each agreement-scale score (1 to 5) to its label and interpretation.
    \begin{tabular}{@{}cll@{}}
      \toprule
      \textbf{Score} & \textbf{Label} & \textbf{Interpretation} \\
      \midrule
      1 & Strongly disagree & Strong negative perception \\
      2 & Disagree & Negative perception \\
      3 & Neither agree nor disagree & Neutral/uncertain \\
      4 & Agree & Positive perception \\
      5 & Strongly agree & Strong positive perception \\
      \bottomrule
    \end{tabular}
  \end{table}

  \subsection{Experience Level Categories (Q6.4)}

  \begin{table}[H]
    \centering
    \caption{Experience level categories.}
    \label{tab:experience-levels}
    % ALT: Maps each experience-range category to its label and career stage.
    \begin{tabular}{@{}lll@{}}
      \toprule
      \textbf{Category} & \textbf{Label} & \textbf{Career Stage} \\
      \midrule
      1--5 Years & Entry Level & Early Career \\
      6--10 Years & Mid-Level (Early) & Mid Career \\
      11--15 Years & Mid-Level (Late) & Mid Career \\
      16+ Years & Senior Level & Late Career \\
      \bottomrule
    \end{tabular}
  \end{table}
